\documentclass[letterpaper]{article}%
\usepackage{lastpage}%
\usepackage{parskip}%
\usepackage{geometry}%
\geometry{left=22.5mm,right=16.1mm,top=48mm,bottom=25mm,headheight=12.5mm,footskip=12.5mm}%
\usepackage{fontspec}%
\usepackage[spanish,activeacute]{babel}%
\usepackage{anyfontsize}%
\usepackage{fancyhdr}%
\usepackage{csquotes}%
\usepackage[ampersand]{easylist}%
\usepackage[style=ieee,backend=biber]{biblatex}%
\usepackage[skins,breakable]{tcolorbox}%
\usepackage{booktabs}%
\usepackage{color}%
\usepackage{tikz}%
\usepackage{tabularx}%
\usepackage{ragged2e}%
\usepackage{graphicx}%
\usepackage{xcolor}%
%
\setmainfont{Arial}%
\addbibresource{../bibliografia.bib}%
\renewcommand*{\bibfont}{\fontsize{10}{14}\selectfont}%

\defbibenvironment{bibliography}
    {\list
    {\printfield[labelnumberwidth]{labelnumber}}
    {\setlength{\leftmargin}{4cm}
    \setlength{\rightmargin}{1.1cm}
    \setlength{\itemindent}{0pt}
    \setlength{\itemsep}{\bibitemsep}
    \setlength{\parsep}{\bibparsep}}}
    {\endlist}
{\item}
%

\newenvironment{textoMargen}
    {%
    \begin{list}{}{%
        \setlength{\leftmargin}{3.6cm}%
        \setlength{\rightmargin}{1.1cm}%
    }%
    \item[]%
  }
  {%
    \end{list}%
  }
%
\definecolor{gris}{rgb}{0.27,0.27,0.27}%
\definecolor{parte}{rgb}{0.02,0.204,0.404}%
\definecolor{azulsuaveTEC}{rgb}{0.02,0.455,0.773}%
\definecolor{fila}{rgb}{0.929,0.929,0.929}%
\definecolor{linea}{rgb}{0.749,0.749,0.749}%
\fancypagestyle{headfoot}{%
\renewcommand{\headrulewidth}{0pt}%
\renewcommand{\footrulewidth}{0pt}%
\fancyhead{%
}%
\fancyfoot{%
}%
\fancyhead[L]{%
\begin{minipage}{0.5\textwidth}%
\flushleft%
\includegraphics[width=62.5mm]{../figuras/Logo.png}%
\end{minipage}%
}%
\fancyfoot[R]{%
\textcolor{black}{%
Página \thepage \hspace{1pt} de \pageref*{LastPage}%
}%
}%
}%
%
\begin{document}%
\normalsize%
\thispagestyle{empty}%
\begin{tikzpicture}[overlay,remember picture]%
\node[inner sep = 0mm,outer sep = 0mm,anchor = north west,xshift = -23mm,yshift = 22mm] at (0.0,0.0) {\includegraphics[width=21cm]{../figuras/Logo_portada.png}};%
\end{tikzpicture}%
\vspace*{100mm}%
\par\fontsize{14}{0}\selectfont \textcolor{black}{Programa del curso EE{-}0405}%
\par\fontsize{18}{25}\selectfont \textbf{\textcolor{black}{Instrumentación}}%
\vspace*{15mm}%
\newline%
\begin{tabularx}{\textwidth}{m{0.02\textwidth}m{0.98\textwidth}}%
&\hspace*{0mm}\fontsize{12}{0}\selectfont \textbf{\textcolor{gris}{Escuela de Ingeniería Electromecánica}}\\%
[-12pt]%
&\hspace*{0mm}\fontsize{12}{0}\selectfont \textbf{\textcolor{gris}{Carrera de Ingeniería Electromecánica}}\\%
\end{tabularx}%
\newpage%
\pagestyle{headfoot}%
\par\fontsize{14}{0}\selectfont \textbf{\textcolor{parte}{I parte: Aspectos relativos al plan de estudios}}%
\par\hspace*{2mm}\fontsize{12}{14}\selectfont \textbf{\textcolor{parte}{1. Datos generales}}%
\vspace*{3mm}%
\newline%
\fontsize{10}{12}\selectfont %
\begin{tabularx}{\textwidth}{p{6cm}p{10cm}}%
\textbf{Nombre del curso:}&Instrumentación\\%
[10pt]%
\textbf{Código:}&EE{-}0405\\%
[10pt]%
\textbf{Tipo de curso:}&Teórico {-} Práctico\\%
[10pt]%
\textbf{Obligatorio o electivo:}&Obligatorio\\%
[10pt]%
\textbf{Nº de créditos:}&2\\%
[10pt]%
\textbf{Nº horas de clase por semana:}&3\\%
[10pt]%
\textbf{Nº horas extraclase por semana:}&3\\%
[10pt]%
\textbf{Ubicación en el plan de estudios:}&Curso de 4\textsuperscript{to} semestre en Ingeniería Electromecánica\\%
[10pt]%
\textbf{Requisitos:}&EE{-}0305 Transductores\\%
[10pt]%
\textbf{Correquisitos:}&Ninguno\\%
[10pt]%
\textbf{El curso es requisito de:}&EE{-}0305 Modelado y simulación de sistemas\\%
[10pt]%
\textbf{Asistencia:}&Obligatoria\\%
[10pt]%
\textbf{Suficiencia:}&No\\%
[10pt]%
\textbf{Posibilidad de reconocimiento:}&Si\\%
[10pt]%
\textbf{Aprobación y actualización del programa:}&01/01/2026 en sesión de Consejo de Escuela 01{-}2026\\%
[10pt]%
\end{tabularx}%
\newpage%
\begin{tabularx}{\textwidth}{p{3cm}p{13cm}}%
\par\fontsize{12}{14}\selectfont \textbf{\textcolor{parte}{2. Descripción general}}&El curso de \emph{Instrumentación} colabora en el desarrollo de los siguientes rasgos del plan de estudios: implementar sistemas de instrumentación para la medición y modificación de variables físicas en sistemas electromecánicos; y aplicar principios de metrología para medir variables físicas en sistemas electromecánicos. \newline\newline Los aprendizajes que los estudiantes desarrollarán en el curso son: integrar transductores, sistemas de adquisición de datos e instrumentos virtuales para la implementación de sistemas de instrumentación; automatizar la adquisición, el procesamiento, la visualización y el control de señales usando instrumentos virtuales; y evaluar el error en la medición y modificación de variables físicas usando los principios de la metrología. \newline\newline Para desempeñarse adecuadamente en este curso, los estudiantes deben poner en práctica lo aprendido en los cursos de: Introducción a la computación, y Transductores. \newline\newline Una vez aprobado este curso, los estudiantes podrán emplear algunos de los aprendizajes adquiridos en los cursos de: Modelado y simulación de sistemas, Control automático, y Control por eventos discretos. \\%
\end{tabularx}%
\vspace*{4mm}%
\newline%
\begin{tabularx}{\textwidth}{p{3cm}p{13cm}}%
\par\fontsize{12}{14}\selectfont \textbf{\textcolor{parte}{3. Objetivos}}&Al final del curso la persona estudiante será capaz de:\newline\newline \textbf{Objetivo general}\begin{itemize}\item Implementar sistemas de instrumentación precisos y confiables para la medición y modificación de variables físicas en sistemas electromecánicos\end{itemize} \vspace{2mm}\textbf{Objetivos específicos}\begin{itemize}\item Integrar transductores, sistemas de adquisición de datos e instrumentos virtuales para la implementación de sistemas de instrumentación\item Automatizar la adquisición, el procesamiento, la visualización y el control de señales usando instrumentos virtuales\item Evaluar el error en la medición y modificación de variables físicas usando los principios de la metrología\end{itemize}\\%
\end{tabularx}%
\newpage%
\begin{tabularx}{\textwidth}{p{3cm}p{13cm}}%
\par\fontsize{12}{14}\selectfont \textbf{\textcolor{parte}{4. Contenidos}}&En el curso se desarrollaran los siguientes temas:\\%
\end{tabularx}%
\newline%
\par \setlength{\leftskip}{4cm} \begin{easylist} \ListProperties(Progressive*=3ex)

& Sistemas de adquisición de datos

&& Arquitectura básica

&& Transductores analógicos y digitales

&& Acondicionamiento de señales

&& A/D: muestreo, cuantización y codificación

&& D/A: decodificación y filtrado

&& Interfaz de transductores con DAQs y microcontroladores

& Instrumentos virtuales

&& Estructuras de datos

&& Estructuras de ejecución

&& Máquinas de estados

&& Guardado de archivos

&& Interfaces gráficas

&& Comunicación serial y USB

&& Control de instrumentos externos

& Simbología y metrología de instrumentación

&& Norma ISA 5.1

&& Calibración y trazabilidad

&& Verificación y validación

&& Propagación de incertidumbres

&& Norma GUM

\end{easylist} \setlength{\leftskip}{0cm} %
\newpage%
\par\fontsize{14}{0}\selectfont \textbf{\textcolor{parte}{II parte: Aspectos operativos}}%
\vspace*{4mm}%
\newline%
\fontsize{10}{12}\selectfont %
\begin{tabularx}{\textwidth}{p{3cm}p{13cm}}%
\par\fontsize{12}{14}\selectfont \textbf{\textcolor{parte}{5. Metodología}}&En este curso, se utilizará el enfoque sistémico-complejo para la ejecución de las sesiones magistrales y se integrará la investigación práctica aplicada para las sesiones prácticas. Esta última se implementará mediante técnicas como la experimentación controlada y el estudio de casos.\newline\newline \textbf{Las personas estudiantes podrán desarrollar actividades en las que:} \newline\begin{itemize}\item Recibirán instrucción sobre la integración de diferentes componentes en un sistema de adquisición de datos y sobre los componentes de un instrumento virtual.\item Integrarán instrumentos virtuales, transductores, dispositivos de adquisición de datos y microcontroladores para crear sistemas de instrumentación a la medida según la aplicación vista en los estudios de caso.\item Calibrarán sistemas de instrumentación y determinarán su incertidumbre.\end{itemize}\vspace*{2mm}Este enfoque metodológico permitirá a la persona estudiante implementar sistemas de instrumentación precisos y confiables para la medición y modificación de variables físicas en sistemas electromecánicos\vspace*{2mm} \newline  Si un estudiante requiere apoyos educativos, podrá solicitarlos a través del Departamento de Orientación y Psicología. \newline \\%
\end{tabularx}%
\vspace*{2mm}%
\newline%
\begin{tabularx}{\textwidth}{p{3cm}p{13cm}}%
\par\fontsize{12}{14}\selectfont \textbf{\textcolor{parte}{6. Evaluación}}&La evaluación se distribuye en los siguientes rubros: \newline \begin{itemize} \item Tareas: implementación de instrumentos virtuales simulados, interpretación de simbología de instrumentación, cálculos de incertidumbres. \item Reportes: desarrollo y conclusión de las prácticas relacionadas con los temas de sistemas de instrumentación integrados. \item Defensa: exposición y defensa de una de las prácticas relacionadas con los temas de sistemas de instrumentación integrados. \end{itemize}\\%
\end{tabularx}%
\vspace*{2mm}%
\newline%
 \begin{minipage}{\linewidth}  \centering  \begin{tabular}{ p{4cm}  p{1.5cm} }  \toprule  Tareas & 20 \% \\  \midrule  Reportes & 60 \% \\  \midrule  Defensa & 20 \% \\  \midrule Total & 100 \% \\  \bottomrule  \end{tabular} \end{minipage}%
\vspace*{2mm}%
\newline%
\begin{tabularx}{\textwidth}{p{3cm}p{13cm}}%
&De conformidad con el artículo 78 del Reglamento del Régimen Enseñanza-Aprendizaje del Instituto Tecnológico de Costa Rica y sus Reformas, en este curso la persona estudiante \textbf{no} tiene derecho a presentar un examen de reposición.\\%
\end{tabularx}%
\vspace*{4mm}%
\newpage%
\begin{tabularx}{\textwidth}{p{3cm}p{13cm}}%
\par\fontsize{12}{14}\selectfont \textbf{\textcolor{parte}{7. Bibliografía}}&\nocite{ida2020sensors} \nocite{fraden2016sensors} \nocite{pallas2012sensors} \nocite{karre2023instru} \nocite{ansiisa51} \nocite{gum2008} \\%
\end{tabularx}%
\vspace*{-8mm}\printbibliography[heading=none]%
\begin{tabularx}{\textwidth}{p{3cm}p{13cm}}%
\par\fontsize{12}{14}\selectfont \textbf{\textcolor{parte}{8. Persona docente}}&El curso será impartido por:\\%
\end{tabularx}%
\vspace*{-4mm}\begin{textoMargen}\textbf{Dr.{-}Ing. Luis Diego Murillo Soto} \vspace*{2mm} \newline Máster en Ciencias de la Ingeniería Eléctrica. Universidad de Costa Rica. Costa Rica \newline \newline  Máster en computación. Ingeniero en Mantenimiento Industrial. Tecnológico de Costa Rica.Costa Rica \newline \newline  Técnico en Electrónica. COVAO \newline \newline \emph{Correo:} lmurillo@itcr.ac.cr\emph{  Teléfono:} 0 \vspace*{1mm} \newline \emph{  Oficina:} 7\emph{  Escuela:} Ingeniería Electromecánica\emph{  Sede:} Cartago \vspace*{4mm} \newline \textbf{Dr.{-}Ing. Juan José Rojas Hernández} \vspace*{2mm} \newline Doctor en ciencia aplicada a la integración de sistemas. Instituto Tecnológico de Kyushu. Japón. \newline \newline  Máster en electrónica con énfasis en microsistemas. Licenciado en Mantenimiento Industrial. Instituto Tecnológico de Costa Rica. Costa Rica. \newline \newline \emph{Correo:} juan.rojas@itcr.ac.cr\emph{  Teléfono:} 88581419 \vspace*{1mm} \newline \emph{  Oficina:} 31\emph{  Escuela:} Ingeniería Electromecánica\emph{  Sede:} Cartago \vspace*{4mm} \newline \end{textoMargen}%
\end{document}